\documentclass[11pt]{article}
\usepackage[margin=1in]{geometry}
\usepackage{amsmath,amssymb}
\usepackage{graphicx}
\usepackage{booktabs}
\usepackage{hyperref}
\usepackage{natbib}
\usepackage{setspace}

\title{Neurofeedback Training Promotes Sharp-Wave Ripples and Preserves Task-Relevant Hippocampal Replay in Rats Performing a Spatial Memory Task}
\author{Author A.$^{1}$, Author B.$^{1,2}$, Author C.$^{2}$ \\
\small $^{1}$Department of Neuroscience, University of Example \\
\small $^{2}$Kavli Institute for Fundamental Neuroscience, Example University}
\date{}

\begin{document}
\maketitle

\begin{abstract}
Hippocampal sharp-wave ripples (SWRs) and associated replay of spatial trajectories are critical for memory consolidation and planning. While disruption of SWRs impairs memory, methods to promote endogenous SWRs while preserving their natural content remain limited. Here we developed a neurofeedback paradigm in which real-time detection of hippocampal SWRs (150--250~Hz ripple band) during a spatial memory task triggered delivery of a tone and food reward. Rats in the manipulation cohort learned to increase SWR rates at a neurofeedback port compared to both a within-subject delay control port and an independent control cohort performing the same task without neurofeedback. Critically, SWR events during neurofeedback intervals contained robust, diverse, and task-relevant spatial replay, including increased remote replay of unvisited maze arms compared to delay trials. The neurofeedback effect could not be explained by differences in head velocity, as SWR rate differences persisted after speed-matched controls. Behavioral performance on the spatial memory task was not impaired by neurofeedback training, with comparable accuracy on neurofeedback and delay trials. Post-targeted-interval compensation in SWR dynamics revealed short-timescale homeostatic regulation of SWR generation. These results demonstrate that neurofeedback can promote endogenous hippocampal memory events without prescribing their content, opening a path toward clinically translatable interventions for memory disorders characterized by reduced SWR activity.
\end{abstract}

\section{Introduction}

Hippocampal sharp-wave ripples (SWRs) are brief, high-frequency oscillatory events that occur during quiet wakefulness and slow-wave sleep and are increasingly recognized as a fundamental substrate of memory consolidation and retrieval \citep{buzsaki2015}. During SWRs, hippocampal place cells reactivate in temporally compressed sequences---termed replay---that recapitulate spatial trajectories experienced during behavior \citep{foster2006,davidson2009}. Selective disruption of SWRs impairs spatial learning and memory \citep{girardeau2009,jadhav2012,ego2018}, while longer-duration SWRs have been associated with improved memory performance \citep{fernandezruiz2019}.

These findings have motivated interest in approaches that could enhance SWR activity as a potential therapeutic strategy for memory impairment. Reduced or abnormal SWRs have been observed in rodent models of Alzheimer's disease, epilepsy, and aging, suggesting that boosting SWR rates could provide cognitive benefit. However, existing positive manipulations have significant limitations. Electrical stimulation during detected SWRs can improve memory, but the stimulation itself may alter the content of replay events. Optogenetic approaches require invasive genetic modifications not applicable to humans. A prior study demonstrated operant conditioning of SWRs using intracranial electrical stimulation as the reinforcer \citep{ishikawa2014}, but intracranial reward is not clinically viable.

A critical unresolved question is whether it is possible to increase SWR rates while preserving the natural diversity and task relevance of replay content. If neurofeedback simply triggers stereotyped or random hippocampal activation, the resulting SWR events might lack the informational content necessary for memory function. Here we address this question using a neurofeedback paradigm embedded within a flexible spatial memory task on an 8-arm maze, enabling direct assessment of both SWR rates and replay content during memory-guided behavior.

\section{Methods}

\subsection{Subjects and Recording}

Rats were implanted with tetrode drives targeting hippocampal area CA1 for simultaneous recording of local field potentials (LFPs) and single-unit spiking activity. Two cohorts were studied: a manipulation cohort receiving neurofeedback and a control cohort performing the same task without neurofeedback. The control cohort data were drawn from \citet{gillespie2021}.

\subsection{Spatial Memory Task}

The task was performed on a flexible 8-arm radial maze. Rats initiated trials by poking their nose into a home port, then maintained a nosepoke at a randomly assigned center port (one of two). After receiving a reward at the center port, rats chose one of eight arms to visit, requiring memory-guided search for rewarded locations.

\begin{table}[ht]
\centering
\caption{Task structure and neurofeedback protocol parameters.}
\label{tab:protocol}
\begin{tabular}{ll}
\toprule
Parameter & Value \\
\midrule
Maze type & 8-arm radial \\
Center ports & 2 (NF port + delay port) \\
Trial initiation & Nosepoke at home port \\
NF trigger & Suprathreshold SWR in CA1 LFP \\
Ripple band & 150--250 Hz \\
NF reward & Tone + food, immediate \\
Delay reward & Tone + food, after matched wait \\
\bottomrule
\end{tabular}
\end{table}

\subsection{Neurofeedback Protocol}

In the manipulation cohort, one center port served as the neurofeedback (NF) port and the other as the delay port (Table~\ref{tab:protocol}). At the NF port, the rat maintained its nosepoke until a suprathreshold SWR was detected in real-time from the ripple-filtered (150--250~Hz) CA1 LFP, at which point a tone and food reward were delivered immediately. At the delay port, the rat waited for a duration randomly drawn from recent NF trial durations before receiving the same reward, providing a within-subject control matched for time at the port.

\subsection{SWR Detection and Analysis}

SWRs were detected in real-time by thresholding the envelope of the 150--250~Hz bandpass-filtered LFP signal recorded from CA1. SWR rates were computed in 0.5-second bins aligned to reward delivery time for comparison across conditions.

\subsection{Replay Decoding}

Replay was decoded using a clusterless approach \citep{denovellis2021} applied to spiking activity during detected SWR events. The 2D maze was linearized to a 1D representation, and Bayesian decoding produced posterior probability distributions over maze positions for each SWR event. Replay events were classified as local (representing the rat's current location or recently visited arms) or remote (representing arms not recently visited).

\subsection{Behavioral Performance}

Search accuracy was quantified as the fraction of trials on which the rat chose a previously unsampled arm during the goal search phase, reflecting memory-guided behavior.

\section{Results}

\subsection{Neurofeedback Training Increases SWR Rates}

Rats in the manipulation cohort developed substantially increased SWR rates during the pre-reward period at center ports. The effect was strongest on NF trials, where a suprathreshold SWR was required for reward delivery, but was also present on delay trials within the same sessions, suggesting a generalized upregulation of SWR generation at center ports.

\begin{table}[ht]
\centering
\caption{SWR rate comparisons across conditions. All comparisons for NF vs.\ delay and NF vs.\ control were statistically significant.}
\label{tab:swr_rate}
\begin{tabular}{lll}
\toprule
Comparison & Direction & Significance \\
\midrule
NF vs.\ delay (manipulation) & NF $\gg$ delay & $p \ll 0.001$ \\
NF vs.\ control & NF $\gg$ control & $p = 1.26 \times 10^{-7}$ \\
Delay vs.\ control & delay $>$ control & $p = 1.12 \times 10^{-4}$ \\
\bottomrule
\end{tabular}
\end{table}

Control cohort animals, which received delay-matched rewards at both center ports without any SWR contingency, showed very few SWRs during the pre-reward period (Table~\ref{tab:swr_rate}). The three-way comparison confirmed that the NF manipulation produced a robust, graded enhancement of SWR rates.

\subsection{Speed Does Not Account for SWR Rate Differences}

Head velocity differed between NF and delay trials, raising the possibility that immobility differences might explain SWR rate differences. However, after matching speed distributions between conditions, NF trials still showed significantly higher SWR rates than delay trials and control trials. This ruled out a purely motoric explanation for the neurofeedback effect.

\begin{table}[ht]
\centering
\caption{Head velocity comparisons between conditions.}
\label{tab:speed}
\begin{tabular}{lll}
\toprule
Comparison & $p$-value & Notes \\
\midrule
NF vs.\ delay (trial-level) & $2.4 \times 10^{-223}$ to $9.1 \times 10^{-36}$ & Multiple time bins \\
NF vs.\ control (groupwise) & $1.26 \times 10^{-7}$ & Significant \\
Delay vs.\ control (groupwise) & $1.12 \times 10^{-4}$ & Significant \\
\bottomrule
\end{tabular}
\end{table}

\subsection{Neurofeedback-Evoked SWRs Contain Task-Relevant Replay}

Clusterless decoding of SWR events during the NF interval revealed robust spatial replay content. Individual SWR events encoded compressed representations of linearized maze positions, including both local replay (representing the rat's current arm or recently visited locations) and remote replay (representing arms not recently visited). NF trials showed a higher rate of remote replay compared to delay trials, indicating that the neurofeedback manipulation enhanced reactivation of distal maze representations.

Replay content was diverse rather than stereotyped: across SWR events within single NF intervals, decoded trajectories spanned multiple maze arms and were not confined to a single repeated sequence. This diversity indicates that the neurofeedback paradigm did not prescribe specific replay content but rather promoted the generation of endogenous SWR events whose content reflected the animal's ongoing task-relevant representations.

\subsection{Behavioral Performance Is Preserved}

A critical concern was whether the cognitive demands of neurofeedback might impair performance on the spatial memory task itself. Analysis of search accuracy---the fraction of trials choosing a previously unsampled arm---revealed no significant differences between NF and delay trials within the manipulation cohort, nor between the manipulation and control cohorts.

\begin{table}[ht]
\centering
\caption{Behavioral performance comparisons. No significant differences were observed in search accuracy.}
\label{tab:behavior}
\begin{tabular}{lll}
\toprule
Comparison & $p$-value & Interpretation \\
\midrule
NF vs.\ delay (within-subject) & 0.76--0.96 & Not significant \\
NF vs.\ control (between-group) & 0.54 & Not significant \\
Delay vs.\ control & 0.54 & Not significant \\
\bottomrule
\end{tabular}
\end{table}

These results (Table~\ref{tab:behavior}) demonstrate that the neurofeedback manipulation did not impose a measurable cognitive burden on spatial memory performance.

\subsection{Post-Interval Compensatory Regulation}

Following the targeted neurofeedback interval, changes in SWR dynamics were observed that suggested short-timescale homeostatic regulation. The elevated SWR rate during the NF interval was partially compensated in subsequent behavioral epochs, indicating that the brain actively regulates SWR occurrence and that neurofeedback-driven increases are not cost-free at the systems level.

\section{Discussion}

We have demonstrated that a neurofeedback paradigm using real-time SWR detection and external food reward can substantially increase hippocampal SWR rates during a spatial memory task while preserving the diversity and task relevance of replay content. This represents an advance over prior approaches to SWR manipulation that either disrupted SWRs \citep{girardeau2009,jadhav2012} or used clinically impractical reinforcers \citep{ishikawa2014}.

\subsection{Neurofeedback Versus Other SWR Manipulations}

Our approach differs fundamentally from electrical or optogenetic stimulation of SWRs. Stimulation-based methods can evoke SWR-like events, but these may not engage the natural brain state permissive for replay, potentially producing physiologically abnormal patterns. In contrast, neurofeedback relies on the animal learning to enter a brain state that naturally generates SWRs. The resulting events therefore arise from endogenous hippocampal dynamics and are more likely to contain meaningful replay content---a prediction confirmed by our decoding analysis.

The use of external food reward rather than intracranial stimulation represents an important step toward clinical translatability \citep{sitaram2017}. While human neurofeedback typically targets sustained oscillatory power (e.g., alpha or theta rhythms), our paradigm targets discrete, brief events, which requires faster detection algorithms but may be more directly relevant to memory-related processes.

\subsection{Increased Remote Replay During Neurofeedback}

The observation that NF trials showed increased remote replay relative to delay trials is particularly noteworthy. Remote replay---reactivation of locations not recently visited---has been linked to memory consolidation and planning \citep{carr2011,foster2006}. The increased rate of remote replay during neurofeedback intervals suggests that the manipulation may preferentially promote the memory consolidation component of SWR function.

\subsection{Homeostatic Regulation of SWRs}

The post-interval compensatory changes we observed suggest that SWR generation is subject to homeostatic regulation on short timescales. This finding has implications for therapeutic applications: sustained neurofeedback sessions may need to account for compensatory downregulation. Understanding the mechanisms of this regulation---whether metabolic, synaptic, or neuromodulatory---will be important for optimizing neurofeedback protocols.

\subsection{Limitations and Future Directions}

Our study did not test whether the increased SWR rates during neurofeedback translate to improved memory performance. The lack of a behavioral deficit is encouraging, but demonstrating a positive effect would require tasks where baseline SWR rates are insufficient for optimal performance, such as in aged animals or disease models. Additionally, we cannot exclude the possibility that the neurofeedback effect is partially mediated by arousal or motivational state changes. Future studies could address this by using yoked controls that receive rewards on the same schedule as NF animals but without SWR contingency.

\section{Conclusion}

We have established that neurofeedback training can promote endogenous hippocampal sharp-wave ripples while preserving the diversity and task relevance of associated spatial replay. The paradigm uses clinically translatable external rewards, does not impair ongoing memory-guided behavior, and reveals homeostatic regulation of SWR generation. These findings provide a foundation for developing neurofeedback-based interventions targeting hippocampal memory processes in conditions characterized by reduced SWR activity.

\bibliographystyle{plainnat}
\bibliography{references}

\end{document}
